\chapter{Transaction Costs in Research}\label{ch:rs:transaction-costs-in-research}

A daily forecast of fifty large stocks has a gross Sharpe ratio of 2.56. Traded as the mean--variance book it implies, by a fund of \$1 billion, it turns over
5.33 times its capital a day and pays 1\,223 per cent of its capital a year in costs: a Sharpe ratio of $-29$ after costs. The same forecast with the costs inside
the optimiser nets 1.80; smoothed over a day as well, 2.43. Nothing about the signal changed. Costs are not a haircut applied to a finished \omterm{def:rs:vectorised-backtests:backtest}{backtest}; they belong
in the model of the market, in the optimiser and in the signal, and they depend on the fund's size. This chapter fits a cost model to the firm's fills, puts it
inside the construction with \texttt{firm.tcost}, nets strategies' trades against each other, and makes signals cost-aware.

\section{Cost models}

\begin{definition}[Temporary impact, permanent impact, square-root impact law, percentage of volume]\label{def:rs:transaction-costs-in-research:impact}
The \emph{temporary impact}\index{temporary impact} of a trade is the part of its price concession that decays after the trade; the \emph{permanent
impact}\index{permanent impact} is the part that remains, the information the market reads in it. An order's \emph{percentage of volume}\index{percentage of volume}
is its size $Q$ over the market's daily volume $V$. The \emph{square-root impact law}\index{square-root impact law} states that an order's average impact, as a fraction
of price, is about $\eta\,\sigma\sqrt{Q/V}$ for daily volatility $\sigma$ and a constant $\eta$ of order one.
\end{definition}

A trade's cost has three parts: half the bid--ask spread (Book~1, chapter~1), fees, and \omterm{def:rs:order-book-replay-simulation:impact}{market impact} (chapter~18). The first two are linear in the size traded;
impact is not. Tóth and co-authors explained the square root by a supply and demand that vanish at the current price, so that a large order must be cut up and
digested; Almgren, Thum, Hauptmann and Li, fitting almost 700\,000 orders executed by Citigroup's US equity desks, found the \omterm{def:rs:transaction-costs-in-research:impact}{temporary impact} grew as the 3/5 power of
the trading rate. The exponent near a half matters more than its precise value: a trade twice as large costs more per share, and so the cost of a book is not
proportional to its size.

\section{Estimating costs from your own fills}

The firm's cost model comes from its own parent orders. The chapter simulates a year of them, 5\,000 orders with sizes from 0.01\% to 20\% of daily volume and
daily volatilities from 1\% to 3\%, from a planted law: half a spread of 2 basis points and impact $0.7\,\sigma\sqrt{Q/V}$, plus the price's own move while the order
works. The move is the difficulty. At 1\% of daily volume and 2\% volatility the planted impact is 14 basis points; the price's move over the order's life has a
standard deviation of 63. A single order says almost nothing about its cost, and a cost model is a regression on thousands of orders, whose noise grows with the
order's size.

\texttt{firm.tcost.fit\_impact} (\cref{lst:rs:transaction-costs-in-research:fit}) averages the costs within twenty bins of participation, fits the exponent on the
logarithms of the bin averages, and fits $\eta$ with the exponent held at a half, with a standard error robust to the growing noise. From 5\,000 orders it recovers
$\eta = 0.71$ (standard error 0.08) and an exponent of 0.48 (0.03); from 500 orders, $1.09$ (0.24) and 0.59 (0.07); from 50\,000, $0.67$ (0.03) and 0.49 (0.01)
(\cref{fig:rs:transaction-costs-in-research:fills}). A desk with a few hundred orders a year knows its cost coefficient to within a third; Kyle and Obizhaeva's
invariance hypotheses, supported on more than 400\,000 portfolio-transition orders, are one way to borrow strength across stocks by scaling costs with each stock's
volume and volatility.

\begin{omfigure}
\begin{tikzpicture}
\begin{loglogaxis}[width=10.5cm,height=6cm,xlabel={participation $Q/V$},ylabel={impact at 2\% volatility (bp)},tick label style={font=\small},label style={font=\small},
  xmin=8e-5,xmax=0.25,ymin=0.5,ymax=150,grid=major,grid style={gray!25},legend pos=north west,legend style={font=\footnotesize},table/col sep=comma]
\addplot[only marks,mark=*,mark size=1.6pt,omDef,error bars/.cd,y dir=both,y explicit] table[x=participation,y=cost_bp,y error=se_bp]{figdata/research/27-transaction-costs-in-research/fills.csv};
\addplot[omThm,thick] table[x=participation,y=fit_bp]{figdata/research/27-transaction-costs-in-research/fills.csv};
\legend{bin averages of 5\,000 orders,fitted power law}
\end{loglogaxis}
\end{tikzpicture}

\omcaption{Impact estimated from a simulated year of parent orders (half spread removed, scaled to 2\% daily volatility), with two-standard-error \omterm{def:rs:market-data-for-research:bar}{bars} and the fitted
law. The smallest orders' bins are the noisiest relative to their cost. Data: \texttt{rs\_tcost.fills}.}\label{fig:rs:transaction-costs-in-research:fills}
\end{omfigure}

\section{Costs inside the optimiser}

\begin{definition}[Cost-aware optimisation]\label{def:rs:transaction-costs-in-research:aware}
\emph{Cost-aware optimisation}\index{cost-aware optimisation} maximises $\alpha^\top w - \frac\gamma2 w^\top\Sigma w - C(w - w_0)$, with $C$ the cost model of the trade
from the current book: here $C(d) = \sum_i s_i\lvert d_i\rvert + \eta\sigma_i\lvert d_i\rvert^{3/2}\sqrt{A/V_i}$ for half spreads $s_i$, a fund of size $A$ and daily traded
values $V_i$.
\end{definition}

The $3/2$ power makes the problem a second-order cone programme (Book~4, chapter~23). \texttt{firm.tcost.cost\_aware} solves it by accelerated proximal gradient
instead: the quadratic part is smooth, the cost is separable, and its proximal step has a closed form, a soft threshold by the spread followed by the root of a
quadratic in $\sqrt{\lvert d\rvert}$ (\cref{lst:rs:transaction-costs-in-research:aware}). The chapter's books hold chapter~26's fifty names long and short, daily; the
forecast is the true expected return of the three planted signals plus noise one and a half times as large, new every day, and $\gamma = 72$. The daily traded
values of these names have a median of \$474 million.

\begin{center}
\small
\begin{tabular}{@{}lrrrr@{}}
\toprule
fund size & \$10 million & \$100 million & \$1 billion & \$10 billion\\
\midrule
mean--variance, costs ignored: after costs & $-4.96$ & $-15.91$ & $-29.47$ & $-34.0$\\
costs inside: Sharpe ratio before costs & 3.12 & 3.14 & 3.06 & 2.11\\
\quad after costs & $-0.31$ & 0.06 & 1.80 & 1.76\\
\quad costs a year & 46.1\% & 28.8\% & 7.6\% & 1.2\%\\
\quad daily turnover & 2.269 & 0.814 & 0.148 & 0.018\\
costs inside, forecast smoothed over one day & 1.68 & 2.03 & 2.43 & 1.66\\
costs ignored, forecast smoothed over fifty days & 1.78 & 1.63 & 1.15 & $-0.23$\\
\bottomrule
\end{tabular}
\end{center}

The book that ignores costs trades 5.33 times its capital a day; even a \$10 million fund, which pays little more than the spread, loses. The cost-aware book has
a surprise at the small end: it loses at \$10 million and wins at \$1 billion. With cheap trading it believes its forecast and chases its daily noise, which the
optimiser cannot tell from alpha; at \$1 billion the cost of trading forces it to move slowly, and moving slowly averages the noise away, so that its Sharpe ratio
before costs (3.06) is higher than the naive book's (2.56). Costs in the optimiser protect a book from its forecast's noise only when they are large relative to
that noise; below that, the forecast itself must be made slower.

\section{Netting across strategies}

\begin{definition}[Internal crossing]\label{def:rs:transaction-costs-in-research:crossing}
\emph{Internal crossing}\index{internal crossing} nets the trades of a firm's strategies in the same stock before any reach the market: a buy of one strategy and
a sell of another cross at no market cost.
\end{definition}

Two teams trading the same fifty names from their own forecasts of the same truth, each with its own noise, each smoothed over twenty days and each running
\$1 billion, would pay 8.27\% and 8.73\% of a team's capital a year apart. Netted, their trades cost 16.04\%: a saving of 0.96 points, 5.7\% of the total. The saving
is smaller than the offsetting trades suggest, because netting cuts both ways under a concave impact law. Where the teams trade opposite ways, the netted trade is
smaller and cheaper; where they trade the same way, the combined trade is larger and, cost growing as the $3/2$ power, dearer than the two apart: 2.94 points a year
dearer here. Each team's own cost estimate understates what its trades cost the firm when others trade alongside it, the first appearance of crowding (chapter~28).
\texttt{firm.tcost.netting\_saving} shares the saving out in proportion to each strategy's standalone cost, so that netting never makes a strategy's reported costs
worse.

\section{Cost-aware signals}

\begin{definition}[Signal smoothing, break-even cost]\label{def:rs:transaction-costs-in-research:smoothing}
\emph{Signal smoothing}\index{signal smoothing} replaces a forecast by an exponentially weighted average of its recent values, with a half-life chosen for net
performance. A strategy's \emph{break-even cost}\index{break-even cost} is the cost per unit traded at which its return after costs is zero: its gross return over
its turnover.
\end{definition}

Smoothing trades signal for turnover. It delays the forecast, which loses the part of the alpha that decays fast, and it averages the forecast's noise, which
removes trading that had no reason. \Cref{fig:rs:transaction-costs-in-research:smoothing} shows the net Sharpe ratio against the half-life at \$1 billion. The
book that ignores costs rises from $-29.47$ to a peak of 1.15 at fifty days, and falls beyond it (0.93 at a hundred days, 0.69 at two hundred) as the \omterm{def:rs:measuring-market-making-and-execution:costs}{delay costs} more
signal than it saves in costs. The cost-aware book peaks early, at 2.43 with a one-day half-life, because the optimiser already smooths; long half-lives only cost it
signal (0.94 at two hundred days). The optimal smoothing depends on who does the rest of the smoothing, and chapter~26's \omterm{def:rs:portfolio-construction-ii:aim}{aim portfolio} is the same idea made exact
for signals whose decay rates are known.

\begin{omfigure}
\begin{tikzpicture}
\begin{semilogxaxis}[width=11cm,height=6cm,xlabel={smoothing half-life (days; the first point is no smoothing)},ylabel={Sharpe ratio after costs},
  tick label style={font=\small},label style={font=\small},xmin=0.22,xmax=250,ymin=-4.2,ymax=3,xtick={0.3,1,10,100},xticklabels={none,1,10,100},
  grid=major,grid style={gray!25},legend pos=south east,legend style={font=\footnotesize},table/col sep=comma]
\addplot[black!60,thick,mark=o] table[x=half_life,y=naive]{figdata/research/27-transaction-costs-in-research/smoothing.csv};
\addplot[omDef,very thick,mark=*] table[x=half_life,y=aware]{figdata/research/27-transaction-costs-in-research/smoothing.csv};
\draw[black!50,densely dotted] (axis cs:0.22,0) -- (axis cs:250,0);
\legend{costs ignored,costs inside the optimiser}
\end{semilogxaxis}
\end{tikzpicture}

\omcaption{The net Sharpe ratio of the fifty-name book at \$1 billion against the half-life of the forecast's smoothing; values below $-4$ are drawn at $-4$. Data:
\texttt{rs\_tcost.smoothing\_curve}.}\label{fig:rs:transaction-costs-in-research:smoothing}
\end{omfigure}

\section{Tutorial: from 2.56 to $-29$ and back}\label{tut:rs:transaction-costs-in-research:tut}

\begin{tutorial}
\textbf{Goal.} Fit the impact law to simulated fills, trade the fifty-name forecast with costs ignored and inside the optimiser at four fund sizes, smooth the
forecast, and net two teams' trades. \textbf{End state:} the table, \cref{fig:rs:transaction-costs-in-research:fills,fig:rs:transaction-costs-in-research:smoothing}.
\begin{tutsteps}
\item \textbf{The fit}: binned averages for the exponent, the robust fit for $\eta$.
\omcode{code/firm/tcost/firm_tcost.py}{34}{56}{Estimating the impact law from fills.}\label{lst:rs:transaction-costs-in-research:fit}
\item \textbf{The cost-aware book}: the closed-form proximal step and the accelerated iteration.
\omcode{code/firm/tcost/firm_tcost.py}{67}{96}{Mean--variance with spread and power-law impact inside.}\label{lst:rs:transaction-costs-in-research:aware}
\item \textbf{Run} \texttt{rs\_tcost.estimate(n)}, \texttt{run(book, aum, half\_life)}, \texttt{smoothing\_curve}, \texttt{netting} and \texttt{fig\_tcost.py}.
\end{tutsteps}
\textbf{What to change next.} Make the forecast's noise persistent (an autoregression rather than new each day) and watch the optimal smoothing lengthen; add
linear costs only and see the cost-aware book's no-trade region appear.
\end{tutorial}

\section{Build: the cost toolkit}\label{bld:rs:transaction-costs-in-research:tcost}

\begin{build}
\bfield{Purpose} One cost model, estimated from the firm's fills and used everywhere: in \omterm{def:rs:vectorised-backtests:backtest}{backtests}, in the optimiser, in capacity estimates and in the netting of
strategies.
\bfield{Interface} \texttt{impact\_bp(eta, sigma, participation, exponent)}, \texttt{fit\_impact(cost, sigma, participation, half\_spread)},
\texttt{trade\_cost(dw, aum, sigma, adv, half\_spread, eta, fee)}, \texttt{prox\_cost(v, s, c)}, \texttt{cost\_aware(alpha, Sigma, w0, aum, sigma, adv, half\_spread,
eta, gamma, neutral)}, \texttt{net\_trades}, \texttt{netting\_saving(lists, cost\_fn)}, \texttt{smooth(signal, half\_life)}, \texttt{breakeven\_cost}.
\bfield{Rules} Costs depend on fund size and are always charged at the size the strategy will run; the cost model is re-fitted on the firm's fills with standard
errors; no strategy is judged before costs; netting savings are shared so that no strategy's costs rise because of another.
\bfield{Acceptance tests} \texttt{code/firm/tcost/tests/}: the planted law recovered within its standard errors; the proximal step optimal on a fine grid; the cost of a
trade by hand; the cost-aware book equal to mean--variance without costs, and locally optimal with them; netting by hand; the smoothing's step response; break-even.
\bfield{Stretch} Temporary and \omterm{def:rs:transaction-costs-in-research:impact}{permanent impact} with decay (a propagator); costs per venue and time of day; multi-period cost-aware construction.
\end{build}

\begin{omsources}
\item B.~Tóth et al., ``Anomalous price impact and the critical nature of liquidity in financial markets'', \emph{Physical Review X} 1, 021006, 2011.
\item R.~Almgren, C.~Thum, E.~Hauptmann and H.~Li, ``Direct estimation of equity market impact'', \emph{Risk}, 2005.
\item A.~S.~Kyle and A.~A.~Obizhaeva, ``Market microstructure invariance: empirical hypotheses'', \emph{Econometrica} 84(4), 2016.
\item J.-P.~Bouchaud, J.~Bonart, J.~Donier and M.~Gould, \emph{Trades, Quotes and Prices}, Cambridge University Press, 2018.
\end{omsources}

\section{Exercises}

\begin{exercise}[$\star$]\label{exo:rs:transaction-costs-in-research:1}
What does an order of 5\% of daily volume cost in a stock of 2\% daily volatility, with a half spread of 2 basis points and $\eta = 0.7$?
\end{exercise}

\begin{exercise}[$\star$]\label{exo:rs:transaction-costs-in-research:2}
The naive book earns 49.1\% a year before costs and trades 5.33 times its capital a day. What is its \omterm{def:rs:transaction-costs-in-research:smoothing}{break-even cost} per unit traded?
\end{exercise}

\begin{exercise}[$\star$]\label{exo:rs:transaction-costs-in-research:3}
Two strategies each trade $q$ in the same direction under a cost growing as $q^{3/2}$. By how much does the netted trade cost more than the two apart?
\end{exercise}

\begin{exercise}[$\star\star$]\label{exo:rs:transaction-costs-in-research:4}
Why does a single parent order say almost nothing about the firm's cost coefficient? How many would you want?
\end{exercise}

\begin{exercise}[$\star\star$]\label{exo:rs:transaction-costs-in-research:5}
Why does the cost-aware book lose at \$10 million but win at \$1 billion?
\end{exercise}

\begin{exercise}[$\star\star$]\label{exo:rs:transaction-costs-in-research:6}
Derive the proximal step of $s\lvert x\rvert + c\lvert x\rvert^{3/2}$.
\end{exercise}

\begin{exercise}[$\star\star\star$]\label{exo:rs:transaction-costs-in-research:7}
\emph{Coding.} Run \texttt{rs\_tcost.smoothing\_curve} for both books and explain why their optimal half-lives differ.
\end{exercise}

\begin{exercise}[$\star\star\star$]\label{exo:rs:transaction-costs-in-research:8}
\emph{Find the flaw.} ``Impact follows a square root, so our costs grow with the square root of our assets: doubling the fund costs only 41\% more.''
\end{exercise}

\section{Problem: From 2.56 to $-29$}

\begin{problem}[{Weekend problem --- a signal, costed}]\label{pb:rs:transaction-costs-in-research:1}
The chapter's simulated fills and fifty-name books.

\medskip\noindent\textbf{Part I --- The cost model.}
\begin{enumerate}
\item State the planted cost law and the size of the price noise at 1\% participation.
\item What do 500, 5\,000 and 50\,000 orders recover for $\eta$ and the exponent?
\item What did Almgren and co-authors and Tóth and co-authors find for the exponent?
\item Why is the fit done on bin averages, and why a robust standard error?
\end{enumerate}

\medskip\noindent\textbf{Part II --- The books.}
\begin{enumerate}[resume]
\item What do the naive book's turnover, costs and net Sharpe ratio come to at \$1 billion?
\item Give the cost-aware book's net Sharpe ratios at the four fund sizes.
\item Why is its Sharpe ratio before costs higher than the naive book's?
\item What happens at \$10 billion?
\end{enumerate}

\medskip\noindent\textbf{Part III --- Netting and smoothing.}
\begin{enumerate}[resume]
\item What do the two teams' costs come to apart and netted?
\item Why does netting cost more on same-direction trades?
\item What are the net Sharpe ratios against the half-life for both books?
\end{enumerate}

\medskip\noindent\textbf{Part IV --- The verdict.}
\begin{enumerate}[resume]
\item State the \emph{named result}: the net Sharpe ratio against the smoothing half-life, and the optimal smoothing.
\item Which book would you run at \$1 billion?
\item What fund size would you advise for the smoothed naive book?
\item How should a research report present a strategy's costs?
\item Where would a propagator model of impact change the conclusions?
\item How would you estimate the firm's cost model if it had no fills yet?
\item What does netting imply for the capacity of several strategies trading the same names?
\item How does chapter~26's \omterm{def:rs:portfolio-construction-ii:aim}{aim portfolio} relate to smoothing?
\item In one sentence: where do costs belong in research?
\end{enumerate}
\end{problem}

\section{Interview questions}

\begin{interviewq}[$\star$ \iqroles{trader, researcher}]\label{iq:rs:transaction-costs-in-research:1}
How does \omterm{def:rs:order-book-replay-simulation:impact}{market impact} scale with order size? What evidence is there?
\end{interviewq}

\begin{interviewq}[$\star\star$ \iqroles{researcher}]\label{iq:rs:transaction-costs-in-research:2}
Your signal has a daily Sharpe ratio of 2.5 before costs. What do you do before you believe it?
\end{interviewq}

\begin{interviewq}[$\star\star$ \iqroles{trader}]\label{iq:rs:transaction-costs-in-research:3}
How would you estimate your desk's \omterm{def:rs:order-book-replay-simulation:impact}{market impact} from its own fills?
\end{interviewq}

\begin{interviewq}[$\star\star$ \iqroles{researcher, developer}]\label{iq:rs:transaction-costs-in-research:4}
How do you put a $3/2$-power cost into a portfolio optimiser?
\end{interviewq}

\begin{interviewq}[$\star\star$ \iqroles{researcher, trader}]\label{iq:rs:transaction-costs-in-research:5}
Two of the firm's strategies trade the same stocks. What should the firm do about it?
\end{interviewq}

\begin{interviewq}[$\star\star\star$ \iqroles{researcher}]\label{iq:rs:transaction-costs-in-research:6}
Show that with costs growing as the $3/2$ power of the trade, a strategy's return after costs is maximised at a finite fund size, and find it for a constant gross
return.
\end{interviewq}
